% Readaptado de Semana 09 Clase 02 (2025-A): dimension y construccion de bases.
\documentclass[compacto,10pt]{pubciencias-notas}
\usepackage{comandosEPN}
\autor{Departamento de Formación Básica}
\materia{Álgebra Lineal}
\nota{Semana 10 Clase 02}
\fecha{Semestre 2026-B}
\logouno[1.1cm]{Logo_EPN.eps}
\logodos[1.7cm]{Logo_DFB.eps}
\definecolor{colordef}{cmyk}{0.81,0.62,0.00,0.22}
\raggedbottom

\begin{document}
\encabezado
\vspace*{-8mm}
\setcounter{section}{3}
\section{Espacios vectoriales}
\subsection{Cardinalidad de las bases y dimensión}

Un espacio vectorial puede admitir distintas bases. Por ejemplo,
$\{1,t\}$ y $\{1+t,1-t\}$ son bases de $\Pol_1$, ambas con dos elementos.
En esta clase se supone que $V$ es un espacio vectorial real de dimensión
finita; es decir, existe un conjunto finito $S\subseteq V$ que genera $V$.
Al eliminar los vectores redundantes de un conjunto generador finito,
se obtiene una base finita de $V$.

\begin{defi}[Cardinalidad]
    Sea $S$ un conjunto finito. Se denomina \textbf{cardinalidad de $S$}
    al número de elementos que lo componen y se denota por $\card(S)$.
    Por ejemplo, $\card(\{1,t\})=2$.
\end{defi}

El siguiente resultado establece que todas las bases
de $V$ tienen la misma cardinalidad. Se utilizará la abreviatura
\textbf{l.i.} para ``linealmente independiente''.

\begin{teo}[Número de elementos de una base]\label{sem10cl02:cardinalidad}
    Sea $B$ una base finita de $V$. Si $L\subseteq V$ es un conjunto
    finito linealmente independiente y $G\subseteq V$ es un conjunto finito que genera $V$,
    entonces
    \[
        \card(L)\le\card(B)\le\card(G).
    \]
    En particular, todas las bases de $V$ tienen el mismo número de elementos.
\end{teo}

\begin{defi}[Dimensión]\label{sem10cl02:dimension}
    Sea $V$ un espacio vectorial real de dimensión finita y sea $B$ una base de $V$.
    Se define la \textbf{dimensión de $V$}, denotada $\dim(V)$, como
    la cardinalidad de $B$:
    \[
        \dim(V)=\card(B).
    \]
    El resultado anterior garantiza que este número no depende de la
    base elegida. Para el espacio $\{0_V\}$, se toma como base el conjunto
    vacío y se define $\dim(\{0_V\})=0$.
\end{defi}

\begin{obs}
    El conjunto $\{0_V\}$ es linealmente dependiente y, por tanto,
    no constituye una base. La base del espacio nulo es $\emptyset$,
    cuya cardinalidad es cero.
\end{obs}

\begin{ejem}[Un plano dentro de un espacio de dimensión 3]
    Para el subespacio $P=\{(x,y,z)\in\R^3:x+y-z=0\}$, se ha demostrado que
    \[
        B=\{(1,0,1),(0,1,1)\}
    \]
    es una base. Por tanto, $\dim(P)=2$. Sus vectores tienen tres
    componentes, sujetas a la restricción $z=x+y$; en consecuencia,
    la tercera componente queda determinada por las dos primeras.

    Por su parte, $\dim(\R^3)=3$. Así, el número de componentes de un
    vector no determina la dimensión del subespacio al que pertenece.
\end{ejem}

\subsection{Dimensiones de los espacios conocidos}

\begin{teo}[Dimensiones de los espacios usuales]\label{sem10cl02:dimensiones}
    Sean $m,n\ge1$ enteros y sea $r\ge0$ un entero. Con las operaciones
    usuales y escalares reales, se tiene:
    \[
        \begin{aligned}
            \dim(\R^n)&=n,\\
            \dim(\Mat[\R]{m}{n})&=mn,\\
            \dim(\Pol_r)&=r+1.
        \end{aligned}
    \]
\end{teo}

Estas igualdades se obtienen de las cardinalidades de las bases canónicas.
La base de $\Pol_r$ contiene los $r+1$ polinomios $1,t,\ldots,t^r$;
la base de $\Mat[\R]{m}{n}$ contiene una matriz $E_{ij}$ por cada
una de las $mn$ posiciones.

\begin{ejem}[Dos bases de un mismo subespacio de matrices]
    Sea $D$ el subespacio de matrices diagonales de $\Mat[\R]{2}{2}$.
    Una base es
    \[
        B=\left\{\begin{pmatrix}1&0\\0&0\end{pmatrix},
                  \begin{pmatrix}0&0\\0&1\end{pmatrix}\right\}.
    \]
    También lo es el conjunto
    \[
        C=\left\{\begin{pmatrix}1&0\\0&1\end{pmatrix},
                  \begin{pmatrix}1&0\\0&-1\end{pmatrix}\right\}.
    \]
    En efecto, para una matriz diagonal genérica,
    \[
        \begin{pmatrix}x&0\\0&y\end{pmatrix}
        =\frac{x+y}{2}\begin{pmatrix}1&0\\0&1\end{pmatrix}
        +\frac{x-y}{2}\begin{pmatrix}1&0\\0&-1\end{pmatrix}.
    \]
    Los coeficientes son únicos, por comparación de las entradas
    diagonales. Ambas bases tienen dos elementos: $\dim(D)=2$.
    En cambio, $\dim(\Mat[\R]{2}{2})=4$.
\end{ejem}

\begin{ejem}[Base y dimensión mediante una parametrización]
    Determine una base y la dimensión de
    \[
        W=\{(a,b,c,d)\in\R^4:a+2b=0,\ c-d=0\}.
    \]
    \textbf{Solución.} Las ecuaciones que definen $W$ implican
    $a=-2b$ y $c=d$. Por tanto,
    \[
        (a,b,c,d)=(-2b,b,d,d)
        =b(-2,1,0,0)+d(0,0,1,1).
    \]
    Los dos vectores generan $W$. La combinación nula
    $\alpha(-2,1,0,0)+\beta(0,0,1,1)=(0,0,0,0)$ implica
    $\alpha=0$ por la segunda componente y $\beta=0$ por la cuarta.
    Así, constituyen una base y $\dim(W)=2$.
\end{ejem}

\begin{advertencia}
    El número de parámetros de una expresión no determina, por sí solo,
    la dimensión. Es necesario verificar que los vectores que multiplican
    a dichos parámetros constituyen una base. Por ejemplo, la expresión
    $(s+t)(1,2)$ contiene dos parámetros, pero describe una recta.
\end{advertencia}

\subsection{Criterios de base en dimensión finita}

\begin{teo}\label{sem10cl02:criterio}
    Sea $\dim(V)=n$ y sea $B\subseteq V$ un conjunto con exactamente
    $n$ elementos. Entonces:
    \begin{enumerate}
        \item si $B$ es linealmente independiente, es una base de $V$;
        \item si $B$ genera $V$, es una base de $V$.
    \end{enumerate}
\end{teo}

Si la dimensión del espacio es conocida y el conjunto tiene esa misma
cardinalidad, basta verificar generación o independencia lineal.
Todo conjunto con más de $n$ elementos es linealmente dependiente, y ningún conjunto con
menos de $n$ elementos genera $V$.

\begin{ejem}[Comprobar una base de polinomios]
    En $\Pol_2$, sea $B=\{1,1+t,1+t+t^2\}$. Como $\dim(\Pol_2)=3$
    y $B$ tiene tres elementos, basta comprobar su independencia lineal. Planteamos
    \[
        \alpha\cdot1+\beta(1+t)+\gamma(1+t+t^2)=0_{\Pol_2}.
    \]
    Al reunir términos e igualar coeficientes obtenemos
    \[
        \left\{\begin{aligned}
            \alpha+\beta+\gamma&=0,\\
            \beta+\gamma&=0,\\
            \gamma&=0.
        \end{aligned}\right.
    \]
    Por sustitución regresiva se obtiene $\gamma=\beta=\alpha=0$. Por tanto, $B$ es
    linealmente independiente y es una base de $\Pol_2$.
\end{ejem}

\subsection{Extracción y extensión de bases}

\begin{prop}
    En un espacio vectorial de dimensión finita:
    \begin{enumerate}
        \item de un conjunto finito que genera el espacio se puede
              extraer una base, eliminando vectores redundantes;
        \item un conjunto linealmente independiente puede completarse hasta una base, añadiendo
              vectores que no pertenezcan a la cápsula de los ya elegidos.
    \end{enumerate}
\end{prop}

\begin{ejem}[Extraer una base de un conjunto generador]
    Sean $a=(1,0,1)$, $b=(0,1,1)$, $c=(1,1,2)$ y $d=(2,-1,1)$.
    Determine una base de $W=\spn\{a,b,c,d\}$.

    \textbf{Solución.} Examinamos los vectores en el orden dado.
    Conservamos cada vector que no pertenece a la cápsula de los
    previamente conservados y descartamos los que sí pertenecen a ella.
    \begin{enumerate}
        \item Como $a=(1,0,1)\ne(0,0,0)$, conservamos $a$.
        \item Para determinar si $b$ es combinación de $a$, buscamos
              $\lambda\in\R$ tal que
              \[
                  \lambda(1,0,1)=(0,1,1).
              \]
              La segunda componente exigiría $0=1$, por lo que no existe
              solución. Así, $b\notin\spn\{a\}$ y conservamos $b$.
        \item Para $c$, planteamos $\alpha a+\beta b=c$. Al igualar
              componentes se obtiene
              \[
                  \left\{\begin{aligned}
                      \alpha&=1,\\\beta&=1,\\\alpha+\beta&=2.
                  \end{aligned}\right.
              \]
              La solución es $\alpha=\beta=1$. Por tanto, $c=a+b$ y
              descartamos $c$, pues no modifica la cápsula.
        \item Para $d$, basta considerar los vectores conservados $a,b$,
              ya que $c\in\spn\{a,b\}$. Planteamos
              $\alpha a+\beta b=d$ y obtenemos
              \[
                  \left\{\begin{aligned}
                      \alpha&=2,\\\beta&=-1,\\\alpha+\beta&=1.
                  \end{aligned}\right.
              \]
              La solución es $\alpha=2$, $\beta=-1$. Así, $d=2a-b$ y
              también descartamos $d$.
    \end{enumerate}
    Como $c=a+b$ y $d=2a-b$, ambos vectores son redundantes. Por la
    Proposición 2 de la semana 07, clase 02,
    su eliminación no modifica el subespacio generado.
    En consecuencia, $W=\spn\{a,b\}$.

    Para verificar que $B=\{a,b\}$ es linealmente independiente,
    planteamos la combinación nula
    \[
        \alpha a+\beta b
        =(\alpha,\beta,\alpha+\beta)=(0,0,0).
    \]
    Al igualar las componentes obtenemos
    \[
        \left\{\begin{aligned}
            \alpha&=0,\\
            \beta&=0,\\
            \alpha+\beta&=0.
        \end{aligned}\right.
    \]
    La única solución es $\alpha=\beta=0$, por lo que $B$ es
    linealmente independiente. Como además genera $W$, concluimos que
    \[
        B=\{(1,0,1),(0,1,1)\}\text{ es una base de }W,
        \qquad\dim(W)=2.
    \]
    Geométricamente, los cuatro vectores pertenecen al plano generado
    por $a,b$; los vectores $c,d$ son redundantes.
\end{ejem}

\begin{ejem}[Extensión de un conjunto linealmente independiente]\label{sem10cl02:completar}
    Sean $u=(1,0,1)$ y $v=(1,-1,0)$. Complete $\{u,v\}$ hasta obtener
    una base de $\R^3$.

    \textbf{Solución.} El conjunto es linealmente independiente, ya que de
    $\alpha u+\beta v=(\alpha+\beta,-\beta,\alpha)=(0,0,0)$ se obtiene
    $\alpha=\beta=0$.
    Para completar el conjunto, consideraremos los vectores de la base
    canónica de $\R^3$:
    \[
        e_1=(1,0,0),\qquad e_2=(0,1,0),\qquad e_3=(0,0,1).
    \]
    Es suficiente añadir uno que no sea combinación lineal de $u,v$.
    Comenzamos con $e_1$ y planteamos
    \[
        su+tv=e_1
        \quad\Longleftrightarrow\quad
        (s+t,-t,s)=(1,0,0).
    \]
    Al igualar componentes obtenemos
    \[
        \left\{\begin{aligned}
            s+t&=1,\\
            -t&=0,\\
            s&=0.
        \end{aligned}\right.
    \]
    Las dos últimas ecuaciones implican $s=t=0$, lo que contradice
    la primera. El sistema no tiene solución; por tanto,
    $e_1\notin\spn\{u,v\}$ y lo incorporamos al conjunto.
    Si $e_1$ hubiese sido combinación de $u,v$, se habría continuado
    con $e_2$ y, de ser necesario, con $e_3$.
    El conjunto $\{u,v,e_1\}$ es linealmente independiente.
    Como el conjunto tiene tres elementos y
    $\dim(\R^3)=3$, entonces  $\{u,v,e_1\}$ constituye una base de $\R^3$.
\end{ejem}

\section{Ejercicios}

\begin{ejer}[Base y dimensión en un subespacio]
    Sea $W=\{(x,y,z)\in\R^3:x-2y+z=0\}$.
    \begin{enumerate}
        \item Determine una base de $W$ y deduzca su dimensión.
        \item Complete esa base hasta obtener una base de $\R^3$.
              Justifique la elección del vector añadido.
    \end{enumerate}
\end{ejer}

\begin{ejer}[Dos bases de un subespacio de matrices]
    Sea
    \[
        W=\left\{\begin{pmatrix}a&b\\c&d\end{pmatrix}
        \in\Mat[\R]{2}{2}:a-2b-c=0,\ a+3d=0\right\}.
    \]
    \begin{enumerate}
        \item Exprese una matriz genérica usando parámetros, obtenga una
              base y determine $\dim(W)$.
        \item Construya otra base distinta y compruebe que ambas tienen
              el mismo número de elementos.
    \end{enumerate}
\end{ejer}

\begin{ejer}[Completar una base de polinomios]
    Sea $T=\{p\in\Pol_2:p(0)=p(1)=0\}$.
    \begin{enumerate}
        \item Escriba $p(t)=a_0+a_1t+a_2t^2$, resuelva las condiciones
              y obtenga una base de $T$ y su dimensión.
        \item Añada polinomios hasta completar esa base a una base de
              $\Pol_2$. Verifique la independencia lineal y utilice la
              dimensión para concluir.
    \end{enumerate}
\end{ejer}

\begin{ejer}[Consecuencias de la dimensión]
    \begin{enumerate}
        \item ¿Puede un conjunto de tres matrices generar
              $\Mat[\R]{2}{2}$? Justifique sin plantear un sistema.
        \item ¿Es necesariamente una base de $\Pol_2$ todo conjunto de
              tres polinomios? Indique qué condición falta y dé un contraejemplo.
    \end{enumerate}
\end{ejer}

\end{document}
